\chapter{Glosario de terminos y abreviaciones}

\section{Abreviaciones}

\begin{longtable}{lp{10cm}}
\toprule
\textbf{Abrev.} & \textbf{Significado} \\
\midrule
\endhead
\bottomrule
\endfoot
CEN     & Coordinador Electrico Nacional (Chile) \\
CNE     & Comision Nacional de Energia (Chile) \\
SEC     & Superintendencia de Electricidad y Combustibles (Chile) \\
SERNAC  & Servicio Nacional del Consumidor (Chile) \\
NERC CIP& Normativa de ciberseguridad electrica (EE.UU.) \\
SEN     & Sistema Electrico Nacional (Chile) \\
SAIDI   & System Average Interruption Duration Index \\
SAIFI  & System Average Interruption Frequency Index \\
SCADA   & Supervisory Control and Data Acquisition \\
PMU     & Phasor Measurement Unit \\
ROCOF   & Rate of Change of Frequency \\
THD     & Total Harmonic Distortion \\
LSTM    & Long Short-Term Memory (red neuronal recurrente) \\
FedAvg  & Federated Averaging (algoritmo de aprendizaje federado) \\
NTIC    & Norma Tecnica de Calidad de Servicio \\
DP      & Diferentially Private (privacidad diferencial) \\
FPR     & False Positive Rate \\
TPR     & True Positive Rate (recall) \\
F1      & Media armonica de precision y recall \\
AUC     & Area Under the Curve (curva ROC) \\
\bottomrule
\end{longtable}

\section{Terminos tecnicos}

\paragraph{Barra AT (alta tension)}
Conjunto de conductores en paralelo que operan a 110~kV en esta
aplicacion. Es el lado de alta tension del transformador.

\paragraph{Barra BT (baja tension)}
Conductores a 23~kV. Es el lado del que salen los alimentadores
hacia los clientes finales.

\paragraph{Corriente nominal de linea}
Corriente maxima que un alimentador puede transportar en regimen
permanente. En este trabajo se asume 800~A.

\paragraph{Demanda activa / reactiva}
La potencia activa (MW) es la que realiza trabajo util; la
reactiva (MVAR) es la debida a cargas magnetizantes (motores,
condensadores). El factor de potencia es $\cos(\phi) = P / S$.

\paragraph{Estado estacionario}
Regimen en el que las variables electricas son constantes en el
tiempo (o cambian lentamente). El modelo de flujo de potencia
describe este regimen.

\paragraph{Factor de potencia}
Cociente entre potencia activa y potencia aparente:
$\cos(\phi) = P / S$. Valores cercanos a 1 indican uso eficiente
de la red.

\paragraph{Intercambiabilidad}
Propiedad estadistica de una secuencia de datos: cualquier
permutacion es igualmente probable bajo la distribucion conjunta.
Es el supuesto que requiere conformal prediction.

\paragraph{Latencia de deteccion}
Tiempo entre el inicio real de la falla y la primera alerta del
sistema. Es la metrica operacional clave.

\paragraph{Norma Tecnica de Calidad de Servicio (NTIC)}
Normativa chilena que fija los indicadores SAIDI/SAIFI maximos
por zona y los plazos de reposicion.

\paragraph{Posicion de tap}
Posicion del cambiador de tomas del transformador, que ajusta la
relacion de transformacion para regular la tension BT.

\paragraph{Residuo}
Diferencia entre el valor observado por SCADA y el valor
predicho por el gemelo digital. Es la materia prima del detector
no supervisado.

\paragraph{Sincrofasor}
Medicion fasorial sincronizada por GPS. Permite medir el angulo
de la tension con precision de microsegundos.

\paragraph{Tap automatico}
Mecanismo del transformador que ajusta la relacion de
transformacion bajo carga, para mantener la tension BT en el
valor de referencia.

\paragraph{Zona de concesion}
Area geografica donde una empresa distribuidora tiene la
obligacion legal de prestar servicio. En Chile las zonas son
asignadas por la SEC.